\section{从语言覆盖率到真实可用性}

要理解 NLLB 的起点，必须先区分三件经常被混为一谈的事：世界上存在多少书面语言，商业产品当时覆盖多少语言，以及其中多少语言真正达到可用、可靠、安全的翻译质量。课堂用一组极简 slide 把这三层逐步展开。

\subsection{三千种书面语言与个人动机}

第一张世界地图不是在追求一个精确不变的语言总数，而是在建立数量级：全球书面语言远多于主流数字产品覆盖的范围。它还暗示一个更深的问题——“有文字”并不等于“有大量数字文本”，更不等于存在可直接训练的高质量平行语料。

\lecturefigure{slide-02-written-languages.jpg}{全球存在三千种以上书面语言}{Stanford 课堂视频 00:01:24}

\begin{knowledgebox}{读图：语言存在、数字化与可训练是三层条件}
地图上的语言首先需要稳定的语言或变体标识；其次需要可获取的数字文本；再次才可能拥有句子级对齐的双语语料。低资源语言可能在第一层非常活跃，却在后两层几乎不可见。因此，训练数据少不是使用者少、语言价值低或语言本身简单的同义词。
\end{knowledgebox}

\begin{warningbox}{不要把“书面语言数量”当成固定真值}
语言与方言边界、书写系统和标准化方式都包含社会选择。课堂数字用于说明覆盖鸿沟，不应用来给每种语言强行指定唯一代码或否定本地社区对名称、脚本和变体的判断。
\end{warningbox}

\teachervoice{老师强调，多语技术并非从未存在；机器翻译本身就是生成式 AI 最成功、最普及的商业应用之一。问题不是“有没有翻译”，而是这种能力为何仍集中在少数语言。}

\subsection{产品覆盖是时间快照，不是质量证明}

接下来课堂把世界语言数量与当时 Google Translate、Microsoft Translator 的覆盖量放在同一张图上。这个对比的第一用途是呈现数量差距，第二用途则是提醒：产品页上的语言计数只说明接口支持，并没有告诉我们长文本、专业领域、低资源方向或安全场景中的误差分布。

\lecturefigure{slide-03-product-coverage.jpg}{课堂时间点的商业翻译产品语言覆盖}{Stanford 课堂视频 00:02:03}

\begin{knowledgebox}{读图：先比较数量级，再追问质量定义}
图中产品覆盖约为百种量级，而书面语言是千种量级。第一眼应看到数量差；第二眼要追问每个“支持”是否包含双向翻译、是否经过母语者验收、是否对低资源方向同样有效。语言数是 coverage 指标，不是 quality、safety 或 equity 指标。
\end{knowledgebox}

\teachervoice{Fan 在报出产品语言数时主动说，这些统计“可能已经有点过时”。因此讲义保留 2023 年课堂时间边界，而不把 slide 数字写成 2026 年的当前产品规格。}

\subsection{目标不是翻倍，而是安全地翻倍}

本节把前面的数量差距转成工程目标。NLLB 的项目问题被压缩成一句话：如果从约百种语言出发，怎样把覆盖翻倍？但 slide 随后补上的括号才是约束核心——不能只增加语言标签，还要尽量得到高质量、安全、真实可用的翻译。

\lecturefigure{slide-04-double-coverage-safely.jpg}{把语言覆盖翻倍，同时保持高质量与安全}{Stanford 课堂视频 00:02:51}

\begin{importantbox}{覆盖目标的正确写法}
设语言集合为 $\mathcal{L}$，系统可声称覆盖的方向集合为 $\mathcal{D}\subseteq\mathcal{L}\times\mathcal{L}$。真正目标不是只最大化 $|\mathcal{L}|$ 或 $|\mathcal{D}|$，而是在每个方向满足最低质量、安全和可用性约束：
\[
\max |\mathcal{D}|\quad
\text{s.t.}\quad q_d\ge \tau_q,\ s_d\ge \tau_s,\ u_d\ge \tau_u\quad \forall d\in\mathcal{D}.
\]
其中 $q_d$ 表示翻译质量，$s_d$ 表示安全表现，$u_d$ 表示使用者可用性；$\tau$ 是项目定义的最低门槛。课堂没有给出统一阈值，公式用于说明多目标约束，而不是补造治理规则。
\end{importantbox}

\begin{warningbox}{语言覆盖不能按标签计数}
同一种语言可能包含不同脚本、地区变体与领域；同一对语言还分 into-English、out-of-English 和 non-English directions。若只把语言名称加到下拉菜单，系统可能在最需要帮助的方向上仍然不可用。
\end{warningbox}

\subsection{高资源描述的是数据，不是人口}

课堂回顾历史进展时指出，机器翻译长期围绕高资源语言发展，Europarl 等机构平行语料成为早期基础。所谓 high-resource 主要描述可获得的训练与评测数据，而不是语言使用人口。一个使用者众多的语言，仍可能因为数字文本、标准化或平行语料不足而处于低资源状态。

\lecturefigure{slide-05-current-progress.jpg}{机器翻译从高资源语言逐步转向低资源语言}{Stanford 课堂视频 00:04:00}

\begin{knowledgebox}{读图：资源轴与人口轴彼此独立}
左侧历史路径依赖机构化双语文本，右侧新进展依赖多语建模、弱监督和语料挖掘。读图时应把“数据资源量”与“人口规模”分开：资源多通常让建模容易，却不说明社会需求更大；资源少则会同时削弱训练、调参与评测。
\end{knowledgebox}

\begin{knowledgebox}{术语消化：资源与数据}
\begin{tabularx}{\textwidth}{@{}L{0.21\textwidth}L{0.28\textwidth}X@{}}
\toprule
术语 & 解决的问题 & 本讲中的含义 \\
\midrule
High-resource language & 描述数据充足度 & 有较多单语、双语与评测资源，不等于人口最多 \\
Low-resource language & 描述数据稀缺度 & 训练、标准化和人工评测资源不足 \\
Monolingual data & 学习单一语言分布 & 网页、文档等只含一种语言的文本 \\
Bitext / parallel data & 学习跨语言映射 & 语义对齐的源句与目标句对 \\
Translation direction & 明确输入输出方向 & $A\rightarrow B$ 与 $B\rightarrow A$ 是两个任务 \\
\bottomrule
\end{tabularx}
\end{knowledgebox}

\subsection{本章小结}

NLLB 的起点不是“世界有三千种语言，所以训练三千种”，而是把覆盖、质量、安全和可用性共同写入目标。高低资源描述数据条件；产品语言数只是时间快照，不能替代方向级评估。

